% Part: normal-modal-logic
% Chapter: tableaux
% Section: rules-for-K

\documentclass[../../../include/open-logic-section]{subfiles}

\begin{document}

\olfileid[pt]{nml}{tab}{rul}

\olsection{Regras para \Ax{K}}

As regras para os conectivos proposicionais usuais são as mesmas dos
!!{tableau}s proposicionais assinalados usuais, apenas com o acréscimo
de prefixos. Em cada caso, a regra aplicada a uma !!{formula} assinalada
$\sFmla{S}{!A}[\sigma]$ produz novas !!{formula}s, também prefixadas
por~$\sigma$. Isso deve ser intuitivamente claro: p.ex., se
$!A \land !B$ é verdadeira em (um mundo nomeado por)~$\sigma$, então
$!A$ e $!B$ são verdadeiras em~$\sigma$ (e não em algum outro mundo).
Reunimos as regras proposicionais em \olref{tab:prop-rules}.

\begin{table}
  \[\def\arraystretch{3}\begin{array}{|c|c|}
    \hline
    \AxiomC{\sFmla{\True}{\lnot !A}[\sigma]}
    \RightLabel{\TRule{\True}{\lnot}}
    \UnaryInfC{\sFmla{\False}{!A}[\sigma]}
    \DisplayProof
    &
    \AxiomC{\sFmla{\False}{\lnot !A}[\sigma]}
    \RightLabel{\TRule{\False}{\lnot}}
    \UnaryInfC{\sFmla{\True}{!A}[\sigma]}
    \DisplayProof
    \\[1ex]
    \hline
    \AxiomC{\sFmla{\True}{!A \land !B}[\sigma]}
    \RightLabel{\TRule{\True}{\land}}
    \UnaryInfC{\sFmla{\True}{!A}[\sigma]}
    \noLine
    \UnaryInfC{\sFmla{\True}{!B}[\sigma]}
    \DisplayProof
    &
    \AxiomC{\sFmla{\False}{!A \land !B}[\sigma]}
    \RightLabel{\TRule{\False}{\land}}
    \UnaryInfC{$\sFmla{\False}{!A}[\sigma] \quad \mid \quad
      \sFmla{\False}{!B}[\sigma]$}
    \DisplayProof
    \\[2ex]
    \hline
    \AxiomC{\sFmla{\True}{!A \lor !B}[\sigma]}
    \RightLabel{\TRule{\True}{\lor}}
    \UnaryInfC{$\sFmla{\True}{!A}[\sigma] \quad \mid \quad
      \sFmla{\True}{!B}[\sigma]$}
    \DisplayProof
    &
    \AxiomC{\sFmla{\False}{!A \lor !B}[\sigma]}
    \RightLabel{\TRule{\False}{\lor}}
    \UnaryInfC{\sFmla{\False}{!A}[\sigma]}
    \noLine
    \UnaryInfC{\sFmla{\False}{!B}[\sigma]}
    \DisplayProof
    \\[2ex]
    \hline
    \AxiomC{\sFmla{\True}{!A \lif !B}[\sigma]}
    \RightLabel{\TRule{\True}{\lif}}
    \UnaryInfC{$\sFmla{\False}{!A}[\sigma] \quad \mid
      \quad \sFmla{\True}{!B}[\sigma]$}
    \DisplayProof
    &
    \AxiomC{\sFmla{\False}{!A \lif !B}[\sigma]}
    \RightLabel{\TRule{\False}{\lif}}
    \UnaryInfC{\sFmla{\True}{!A}[\sigma]}
    \noLine
    \UnaryInfC{\sFmla{\False}{!B}[\sigma]}
    \DisplayProof
    \\[2ex]
    \hline
  \end{array}\]
  \caption{Regras de !!{tableau} prefixado para os conectivos
    proposicionais}
  \ollabel{tab:prop-rules}
\end{table}

A condição de fechamento é a mesma dos !!{tableau}s comuns, embora
exijamos que coincidam não apenas as !!{formula}s, mas também os
prefixos. Assim, um ramo é fechado se contém tanto
\[
\sFmla{\True}{!A}[\sigma] \quad\text{quanto}\quad \sFmla{\False}{!A}[\sigma]
\]
para algum prefixo $\sigma$ e alguma !!{formula}~$!A$.

As regras para estabelecer hipóteses também são as mesmas dos
!!{tableau}s comuns, exceto que, para hipóteses, usamos sempre o
prefixo~$1$. (Não importa qual prefixo usamos, desde que seja o mesmo
para todas as hipóteses.) Assim, p.ex., dizemos que
\[
!B_1, \dots, !B_n \Proves !A
\]
sse há um tableau fechado para as hipóteses
\[
\sFmla{\True}{!B_1}[1], \dots, \sFmla{\True}{!B_n}[1],
\sFmla{\False}{!A}[1].
\]

Para os operadores modais\iftag{prvBox}{\iftag{prvDiamond}{~$\Box$
    e}{~$\Box$}}{}\iftag{prvDiamond}{~$\Diamond$}{}, o prefixo da
conclusão da regra aplicada a !!a{formula} com prefixo~$\sigma$ é
$\sigma.n$. Contudo, o valor permitido de $n$ depende de o sinal ser
~$\True$ ou~$\False$.

\iftag{prvBox}{A regra $\TRule{\True}{\Box}$ estende um ramo que contém
  $\sFmla{\True}{\Box !A}[\sigma]$ com
  $\sFmla{\True}{!A}[\sigma.n]$.\iftag{prvDiamond}{ De modo semelhante,
    }{}}{}\iftag{prvDiamond}{a regra $\TRule{\False}{\Diamond}$ estende
  um ramo que contém $\sFmla{\False}{\Diamond !A}[\sigma]$ com
  $\sFmla{\False}{!A}[\sigma.n]$.}{}
\iftag{notprvBox,notprvDiamond}{Ela só pode ser aplicada}{Elas só podem
ser aplicadas} a um prefixo~$\sigma.n$ que \emph{já} ocorra no ramo em que a regra é
aplicada. Chamemos tal prefixo de ``usado'' (no ramo).

\iftag{prvBox}{A regra $\TRule{\False}{\Box}$ estende um ramo que contém
  $\sFmla{\False}{\Box !A}[\sigma]$ com
  $\sFmla{\False}{!A}[\sigma.n]$.\iftag{prvDiamond}{ De modo semelhante,
    }{}}{}\iftag{prvDiamond}{a regra $\TRule{\True}{\Diamond}$ estende
  um ramo que contém $\sFmla{\True}{\Diamond !A}[\sigma]$ com
  $\sFmla{\True}{!A}[\sigma.n]$.}{}
\iftag{notprvBox,notprvDiamond}{Essa regra, contudo, só pode ser aplicada}{Essas regras,
contudo, só podem ser aplicadas} a um prefixo~$\sigma.n$ que \emph{ainda não} ocorra
no ramo em que a regra é aplicada. Chamamos tais prefixos de ``novos''
(para o ramo).

As regras são apresentadas em \olref{tab:rules-K}.

\begin{table}
  \begin{center}
    \def\arraystretch{3}\def\fCenter{}
    \begin{tabular}{|c|c|}
    \hline
    \iftag{prvBox}{
      \AxiomC{\sFmla{\True}{\Box !A}[\sigma]}
      \RightLabel{\TRule{\True}{\Box}}
      \UnaryInfC{\sFmla{\True}{!A}[\sigma.n]}
      \DisplayProof
      &
      \AxiomC{\sFmla{\False}{\Box !A}[\sigma]}
      \RightLabel{\TRule{\False}{\Box}}
      \UnaryInfC{\sFmla{\False}{!A}[\sigma.n]}
      \DisplayProof\\
      $\sigma.n$ é usado & $\sigma.n$ é novo
      \\[1ex]
      \hline}{}
    \iftag{prvDiamond}{
      \AxiomC{\sFmla{\True}{\Diamond !A}[\sigma]}
      \RightLabel{\TRule{\True}{\Diamond}}
      \UnaryInfC{\sFmla{\True}{!A}[\sigma.n]}
      \DisplayProof
      &
      \AxiomC{\sFmla{\False}{\Diamond !A}[\sigma]}
      \RightLabel{\TRule{\False}{\Diamond}}
      \UnaryInfC{\sFmla{\False}{!A}[\sigma.n]}
      \DisplayProof\\
      $\sigma.n$ é novo & $\sigma.n$ é usado
      \\[1ex]
      \hline}{}
    \end{tabular}
  \end{center}
  \caption{As regras modais para \Ax{K}.}
  \ollabel{tab:rules-K}
\end{table}

A exigência de que o prefixo para
\iftag{prvBox}{\TRule{\True}{\Box}}{\TRule{\False}{\Diamond}} seja usado
é necessária, pois, do contrário, contaríamos o seguinte como um
!!{tableau} fechado:
\iftag{notprvBox,notprvDiamond}{%
  \iftag{prvBox}{%
  \begin{oltableau}
    [\pFmla{\True}{\Box \formula{A}}{1}, just = \TAss
      [\pFmla{\False}{\lnot\Box\lnot \formula{A}}{1}, just = \TAss
        [\pFmla{\True}{\formula{A}}{1.1}, just = {\TRule{\True}{\Box}[1]}
          [\pFmla{\True}{\Box\lnot\formula{A}}{1},
            just ={\TRule{\False}{\lnot}[2]}
            [\pFmla{\True}{\lnot\formula{A}}{1.1},
              just = {\TRule{\True}{\Box}[4]}
              [\pFmla{\False}{\formula{A}}{1.1},
                  just ={\TRule{\True}{\lnot}[5]}, close]
            ]
          ]
        ]
      ]
    ]
  \end{oltableau}
  }{
  \begin{oltableau}
    [\pFmla{\True}{\lnot\Diamond\lnot \formula{A}}{1}, just = \TAss
      [\pFmla{\False}{\Diamond\formula{A}}{1}, just = \TAss
        [\pFmla{\False}{\formula{A}}{1.1}, just = {\TRule{\False}{\Diamond}[2]}
          [\pFmla{\False}{\Diamond\lnot\formula{A}}{1},
            just ={\TRule{\True}{\lnot}[1]}
            [\pFmla{\False}{\lnot\formula{A}}{1.1},
              just = {\TRule{\False}{\Diamond}[4]}
              [\pFmla{\True}{\formula{A}}{1.1},
                  just ={\TRule{\True}{\lnot}[5]}, close]
            ]
          ]
        ]
      ]
    ]
  \end{oltableau}
}}{
  \begin{oltableau}
    [\pFmla{\True}{\Box \formula{A}}{1}, just = \TAss
      [\pFmla{\False}{\Diamond \formula{A}}{1}, just = \TAss
        [\pFmla{\True}{\formula{A}}{1.1}, just = {\TRule{\True}{\Box}[1]}
          [\pFmla{\False}{\formula{A}}{1.1},
            just = {\TRule{\False}{\Diamond}[2]}, close]
        ]
      ]
    ]
\end{oltableau}}

Mas $\Box \formula{A} \Entails/ \Diamond \formula{A}$; portanto, nosso
sistema de demonstração seria incorreto. Analogamente,
$\Diamond \formula{A} \Entails/ \Box \formula{A}$, mas, sem a restrição
de que o prefixo para
\iftag{prvBox}{\TRule{\False}{\Box}}{\TRule{\True}{\Diamond}} seja novo,
este seria um tableau fechado: \iftag{notprvBox,notprvDiamond}{%
  \iftag{prvBox}{%
  \begin{oltableau}
    [\pFmla{\True}{\lnot\Box\lnot \formula{A}}{1}, just = \TAss
      [\pFmla{\False}{\Box \formula{A}}{1}, just = \TAss
        [\pFmla{\False}{\formula{A}}{1.1}, just = {\TRule{\True}{\Box}[2]}
          [\pFmla{\False}{\Box\lnot\formula{A}}{1},
            just ={\TRule{\True}{\lnot}[1]}
            [\pFmla{\False}{\lnot\formula{A}}{1.1},
              just = {\TRule{\False}{\Box}[4]}
              [\pFmla{\True}{\formula{A}}{1.1},
                  just ={\TRule{\False}{\lnot}[5]}, close]
            ]
          ]
        ]
      ]
    ]
  \end{oltableau}
  }{
  \begin{oltableau}
    [\pFmla{\True}{\Diamond \formula{A}}{1}, just = \TAss
      [\pFmla{\False}{\lnot\Diamond\lnot\formula{A}}{1}, just = \TAss
        [\pFmla{\True}{\formula{A}}{1.1}, just = {\TRule{\True}{\Diamond}[1]}
          [\pFmla{\True}{\Diamond\lnot\formula{A}}{1},
            just ={\TRule{\False}{\lnot}[2]}
            [\pFmla{\True}{\lnot\formula{A}}{1.1},
              just = {\TRule{\True}{\Diamond}[4]}
              [\pFmla{\False}{\formula{A}}{1.1},
                  just ={\TRule{\True}{\lnot}[5]}, close]
            ]
          ]
        ]
      ]
    ]
  \end{oltableau}
}}{
  \begin{oltableau}
    [\pFmla{\True}{\Diamond \formula{A}}{1}, just = \TAss,name=one
      [\pFmla{\False}{\Box \formula{A}}{1}, just = \TAss, name=two
        [\pFmla{\True}{\formula{A}}{1.1}, just = {\TRule{\True}{\Diamond}[1]}
          [\pFmla{\False}{\formula{A}}{1.1},
            just = {\TRule{\False}{\Box}[2]}, close]
        ]
      ]
    ]
  \end{oltableau}}

\end{document}
